
\section*{0.4. Profinite Algebras}
\phantomsection
\label{sga3:VIIB:0.4}
\addcontentsline{toc}{subsection}{\SGARefLink{sga3:VIIB:0.4}{0.4}. Profinite Algebras}

Let \(k\) be a pseudocompact ring. A \emph{topological
\(k\)-algebra} is a commutative topological ring \(A\), together with
a morphism of topological rings \(k\to A\). The algebra \(A\) is called
a \emph{profinite \(k\)-algebra} if its underlying topological
\(k\)-module is pseudocompact.

Let \(\mathfrak l\) be an open \(k\)-submodule of such an \(A\). The
composite map
\[
  \varphi:
  A\times A
  \xrightarrow{\mathrm{mult}}A
  \longrightarrow A/\mathfrak l
\]
is continuous. \SGARefLink{sga3:VIIB:0.3.1}{Lemma 0.3.1} therefore gives an open \(k\)-submodule
\(\mathfrak n\) of \(A\) such that
\(\varphi(A\times\mathfrak n)=0\). Thus \(\mathfrak l\) contains the
open ideal \(A\mathfrak n\), and \(A\) is a pseudocompact ring.

Likewise, let \(M\) be a topological \(A\)-module whose underlying
\(k\)-module is pseudocompact. If \(M'\) is an open \(k\)-submodule of
\(M\), applying \SGARefLink{sga3:VIIB:0.3.1}{Lemma 0.3.1} to
\[
  A\times M
  \xrightarrow{\mathrm{mult}}M
  \longrightarrow M/M'
\]
shows that \(M'\) contains an open \(A\)-submodule. Hence \(M\) is also
a pseudocompact \(A\)-module.\footnote{The editors added the following
lemma.}
Conversely:

\medskip
\noindent\textbf{Lemma.}
\phantomsection
\label{sga3:VIIB:0.4-restriction-lemma}
\textit{Let \(A\) be a profinite \(k\)-algebra and \(M\) a
pseudocompact \(A\)-module. The \(k\)-module \(M|_k\), obtained by
restriction of scalars, is pseudocompact.}

Indeed, every finite-length pseudocompact \(A\)-module is isomorphic
to a quotient
\[
  A^n/L,
\]
where \(L\) is an open submodule of \(A^n\), and is therefore a
pseudocompact \(k\)-module. Since \(M|_k\) is a projective limit of
modules of this kind, it is a pseudocompact \(k\)-module.

\medskip
\noindent\textbf{0.4.1.}
\phantomsection
\label{sga3:VIIB:0.4.1}
If \(A\) and \(B\) are profinite \(k\)-algebras, a \emph{morphism}
from \(A\) to \(B\) is, by definition, a continuous homomorphism of
\(k\)-algebras. We denote the category of profinite \(k\)-algebras by
\[
  \operatorname{Alp}/k.
\]
It plainly has projective limits: the algebra underlying a projective
limit is the projective limit of the underlying algebras, endowed with
the projective-limit topology.

It also has finite direct limits. For example, if
\[
  f:A\longrightarrow B,
  \qquad
  g:A\longrightarrow C
\]
are morphisms of profinite \(k\)-algebras, the amalgamated sum of
\(B\) and \(C\) over \(A\) has
\[
  B\widehat{\otimes}_A C
\]
as its underlying topological \(A\)-module. Here \SGARefLink{sga3:VIIB:0.4}{0.4} makes \(B\) and
\(C\) pseudocompact \(A\)-modules via \(f\) and \(g\), and the
multiplication is characterized by
\[
  (b\widehat{\otimes}c)
  (b'\widehat{\otimes}c')
  =
  bb'\widehat{\otimes}cc'
\]
for \(b,b'\in B\) and \(c,c'\in C\).

\medskip
\noindent\textbf{Definition 0.4.2.}
\phantomsection
\label{sga3:VIIB:0.4.2}
\textit{A profinite \(k\)-algebra \(C\) is said to have
\emph{finite length} if its underlying \(k\)-module has finite length,
and is therefore discrete. We denote by \(\operatorname{Alf}/k\) the
full subcategory of \(\operatorname{Alp}/k\) formed by the
\(k\)-algebras of finite length.}

For every profinite \(k\)-algebra \(A\), let \(h_A\) be the functor
\[
  h_A:\operatorname{Alf}/k\longrightarrow(\mathbf{Sets}),
  \qquad
  C\longmapsto
  \operatorname{Hom}_{\operatorname{Alp}/k}(A,C).
\]
The functor \(h_A\) is left exact, and the canonical projections
\(A\to A/\mathfrak l\), where \(\mathfrak l\) ranges over the open
ideals of \(A\), induce for every \(C\in\operatorname{Alf}/k\) an
isomorphism
\[
  \varinjlim_{\mathfrak l}
  \operatorname{Hom}_{\operatorname{Alf}/k}(A/\mathfrak l,C)
  \xrightarrow{\sim}
  \operatorname{Hom}_{\operatorname{Alp}/k}(A,C),
\]
functorial in \(C\). In other words, \(h_A\) is the direct limit of the
representable functors \(h_{A/\mathfrak l}\):
\begin{equation}
  h_A\simeq\varinjlim_{\mathfrak l}h_{A/\mathfrak l}.
  \tag{\(\ast\)}
  \label{sga3:VIIB:0.4.2-star}
\end{equation}

If \(B\) is another profinite \(k\)-algebra, the general properties of
the bifunctor \(\operatorname{Hom}\), together with
\[
  \operatorname{Hom}(h_C,h_B)=h_B(C)
\]
for \(C\) of finite length, give isomorphisms
\begin{align*}
  \operatorname{Hom}_{\operatorname{Alp}/k}(B,A)
  &\simeq
  \varprojlim_{\mathfrak l}
  \operatorname{Hom}_{\operatorname{Alp}/k}(B,A/\mathfrak l)
  \\
  &\simeq
  \varprojlim_{\mathfrak l}
  \operatorname{Hom}(h_{A/\mathfrak l},h_B)
  \\
  &\simeq
  \operatorname{Hom}
    \left(\varinjlim_{\mathfrak l}h_{A/\mathfrak l},h_B\right).
\end{align*}
Together with \eqref{sga3:VIIB:0.4.2-star}, this proves that the
contravariant functor \(A\mapsto h_A\) is fully faithful. In fact:

\medskip
\noindent\textbf{Proposition.}
\phantomsection
\label{sga3:VIIB:0.4.2-proposition}
\textit{The functor \(A\mapsto h_A\) is an anti-equivalence from
\(\operatorname{Alp}/k\) onto the category of left-exact functors from
\(\operatorname{Alf}/k\) to \((\mathbf{Sets})\).}

By what precedes, it remains only to show that every left-exact functor
\[
  F:\operatorname{Alf}/k\longrightarrow(\mathbf{Sets})
\]
is isomorphic to a functor \(h_A\). The algebra \(A\) may be constructed
as follows (cf.\ TDTE~II, \S~3).

Because \(F\) is left exact, for every finite-length \(k\)-algebra \(C\)
and every \(\xi\in F(C)\), there is a smallest subalgebra \(C'\) of
\(C\) such that \(\xi\) belongs to the image of \(F(C')\) in \(F(C)\).
If \(C'=C\), the pair \((C,\xi)\) is called \emph{minimal}.

The minimal pairs form a category when a morphism from \((C,\xi)\) to
\((D,\eta)\) is defined to be a homomorphism
\(\phi:C\to D\) such that
\[
  (F\phi)(\xi)=\eta.
\]
Such a \(\phi\) is necessarily surjective, and the category of minimal
pairs is left filtered. Moreover, one may restrict to pairs
\((C,\xi)\) for which \(C\) belongs to a set containing one
finite-length \(k\)-algebra of every isomorphism
type.
The functor \((C,\xi)\mapsto C\), from the category of minimal pairs to
the category of profinite \(k\)-algebras, therefore has a projective
limit; take \(A\) to be this limit.

\medskip
\noindent\textbf{Corollary.}
\phantomsection
\label{sga3:VIIB:0.4.2-corollary}
\textit{The category \(\operatorname{Alp}/k\) has arbitrary direct
limits.}

Indeed, the category of left-exact functors from
\(\operatorname{Alf}/k\) to \((\mathbf{Sets})\) has projective limits,
computed pointwise: if \((F_i)\) is a projective system of such
functors, then for every \(C\in\operatorname{Alf}/k\),
\[
  (\varprojlim_iF_i)(C)=\varprojlim_iF_i(C).
\]
\footnote{\textit{Editors' note.} If \((A_i)_{i\in I}\) is a direct system in
\(\operatorname{Alp}/k\), its direct limit in \(\operatorname{Alp}/k\)
is the separated completion of the \(k\)-algebra \(B\), the direct
limit of the underlying \(k\)-algebras, for the topology defined by the
ideals \(I\) such that \(I\cap A_i\) is open in \(A_i\) for every \(i\)
and \(\operatorname{length}_k(B/I)<\infty\). Notice, for example, that
if \(K/k\) is an algebraic extension of fields of infinite degree and
\((k_i)\) is the filtered direct system of finite-degree subextensions,
then the direct limit of the system \((k_i)\) in
\(\operatorname{Alp}/k\) is the zero ring.}

